\documentclass[12pt,a4paper]{article}
\usepackage{../latex/dsdlab-report}

\begin{document}
	\reporttitlepage{۱}{طراحی مدارهای ترکیبی با طرح‌واره}

	\section*{۱.۱ مقدمه و اهداف}
	هدف آزمایش، آشنایی با طراحی طرح‌واره‌ای به کمک رایانه و طراحی، شبیه‌سازی، سنتز و نگاشت مدارهای ترکیبی برای آرایهٔ دروازه‌ای برنامه‌پذیر میدانی است. مطابق دستورکار، مدارهای تشخیص مضرب ۳ و مضرب ۱۱ برای یک ورودی چهاررقمی با رمزگذاری دهدهیِ دودویی‌شده، تنها با دروازه‌های پایه طراحی می‌شوند.

	\section*{۱.۲ طراحی مدار}
	\subsection*{الف) تشخیص مضرب ۳}
	برای عدد چهاررقمی با رمزگذاری دهدهیِ دودویی‌شده به‌صورت
	\[
		N=1000A+100B+10C+D
	\]
	داریم \(10\equiv1\pmod 3\)، در نتیجه
	\[
		N\bmod3=(A+B+C+D)\bmod3.
	\]
	برای ساخت جمع‌کننده‌های مورد نیاز، ابتدا نیم‌جمع‌کننده با روابط
	\[
		S=A\overline{B}+\overline{A}B,\qquad C=AB
	\]
	ساخته شد. برای رعایت محدودیّت استفاده از دروازه‌های پایه، دروازهٔ \lr{XOR}\LTRfootnote{Exclusive OR Gate} با دروازه‌های \lr{AND}\LTRfootnote{AND Gate}، \lr{OR}\LTRfootnote{OR Gate} و \lr{NOT}\LTRfootnote{NOT Gate} پیاده‌سازی شد؛ بنابراین نماد \(\oplus\) تنها صورت فشردهٔ همین تابع منطقی است. سپس دو نیم‌جمع‌کننده و یک دروازهٔ \lr{OR} یک جمع‌کنندهٔ کامل را تشکیل دادند. در این مدار، اگر \(S_1=A\oplus B\) و \(C_1=AB\) باشند، خروجی‌ها از روابط
	\[
		S=A\oplus B\oplus C_{\mathrm{in}},\qquad
		C_{\mathrm{out}}=C_1+S_1C_{\mathrm{in}}
	\]
	به دست می‌آیند. چهار نمونه از جمع‌کنندهٔ کامل نیز با زنجیره‌کردن مسیر نقلی به یک جمع‌کنندهٔ چهاربیتی متّصل شدند.

	\shot{screenshots/Div3/1.png}{پیاده‌سازی نیم‌جمع‌کننده با دروازه‌های پایه}
	\shot{screenshots/Div3/2.png}{ساخت جمع‌کننده‌ی کامل از دو نیم‌جمع‌کننده}
	\shot{screenshots/Div3/3.png}{جمع‌کنندهٔ چهاربیتی با اتّصال آبشاری چهار جمع‌کنندهٔ کامل}
	\shot{screenshots/Div3/4.png}{شبیه‌سازی جمع‌کنندهٔ چهاربیتی و مشاهدهٔ جمع و نقلی}

	پیمانهٔ \lr{rem} با منطق ترکیبی، باقی‌ماندهٔ یک رقم دهدهیِ دودویی‌شده بر ۳ را تولید می‌کند. باقی‌مانده در \(O[1..0]\) قرار می‌گیرد و دو بیت بالاییِ \(O[3..2]\) همواره صفرند. جدول زیر مبنای استخراج معادلات دروازه‌ای این پیمانه است:
	\[
	\begin{array}{c|cccccccccc}
	\text{رقم دهدهی} & 0&1&2&3&4&5&6&7&8&9\\ \hline
	\text{باقی‌مانده بر 3} & 0&1&2&0&1&2&0&1&2&0
	\end{array}
	\]
	ورودی‌های ۱۰ تا ۱۵ نمایش معتبر دهدهیِ دودویی‌شده نیستند و در طراحی و آزمون استفاده نمی‌شوند. در مدار \lr{Div3} باقی‌ماندهٔ چهار رقم ابتدا محاسبه، به کمک سه جمع‌کننده با هم جمع و دوباره بر ۳ تقلیل داده می‌شود. خروجی داخلی \lr{div3} تنها هنگامی یک است که دو بیت باقی‌ماندهٔ نهایی صفر باشند.

	\shot{screenshots/Div3/5.png}{منطق دروازه‌ای پیمانهٔ محاسبهٔ باقی‌مانده بر ۳}
	\shot{screenshots/Div3/7.png}{شبیه‌سازی پیمانهٔ باقی‌مانده برای رقم‌های صفر تا ۹}
	\shot{screenshots/Div3/8.png}{طرح‌وارهٔ سلسله‌مراتبی آشکارساز مضرب ۳ چهاررقمی}

	\subsection*{ب) تشخیص مضرب ۱۱}
	از آنجا که \(10\equiv-1\pmod {11}\)، شرط بخش‌پذیری عدد چهاررقمی برابر است با
	\[
		N\bmod11=(D-C+B-A)\bmod11=0,
	\]
	یا به‌صورت معادل، برابری دو جمع متناوب \(A+C\) و \(B+D\) با در نظر گرفتن حالت‌های هم‌نهشت بر ۱۱.

	اگر \(X=A+C\) و \(Y=B+D\) باشد، با توجّه به \(0\leq X,Y\leq18\)، عدد ورودی بر ۱۱ بخش‌پذیر است اگر و تنها اگر یکی از سه شرط \(X=Y\)، \(X=Y+11\) یا \(Y=X+11\) برقرار باشد. برای انجام این بررسی، مقایسه‌کنندهٔ یک‌بیتی با سه خروجی بزرگ‌تر، برابر و کوچک‌تر ساخته شد. دو نمونه از آن یک مقایسه‌کنندهٔ دوبیتی و دو نمونهٔ مقایسه‌کنندهٔ دوبیتی یک مقایسه‌کنندهٔ چهاربیتی را تشکیل می‌دهند؛ در هر سطح، نتیجهٔ بیت پرارزش بر نتیجهٔ بیت کم‌ارزش اولویت دارد. یک تسهیم‌کنندهٔ\LTRfootnote{Multiplexer} چهاربیتی نیز با چهار تسهیم‌کنندهٔ دروازه‌ای یک‌بیتی پیاده‌سازی شد. رابطهٔ هر بیت این تسهیم‌کننده برابر است با
	\[
		O_i=\overline{\mathrm{Sel}}\,N0_i+\mathrm{Sel}\,N1_i,
		\qquad i=0,1,2,3,
	\]
	و در نتیجه این زیرمدار نیز فقط از دروازه‌های پایه ساخته شده است.

	\shot{screenshots/Div11/1.png}{مقایسه‌کننده‌ی یک‌بیتی با خروجی‌های بزرگ‌تر، برابر و کوچک‌تر}
	\shot{screenshots/Div11/2.png}{ساخت مقایسه‌کننده‌ی دوبیتی از دو مقایسه‌کننده‌ی یک‌بیتی}
	\shot{screenshots/Div11/3.png}{ساخت مقایسه‌کننده‌ی چهاربیتی از دو مقایسه‌کننده‌ی دوبیتی}
	\shot{screenshots/Div11/4.png}{شبیه‌سازی حالت‌های بزرگ‌تر، کوچک‌تر و برابر مقایسه‌کننده}
	\shot{screenshots/Div11/5.png}{تسهیم‌کنندهٔ چهاربیتی ساخته‌شده از دروازه‌های پایه}

	در مدار نهایی \lr{Div11}، ابتدا دو جمع \(X=A+C\) و \(Y=B+D\) محاسبه می‌شوند که هر یک در بازهٔ صفر تا ۱۸ قرار دارد. برای تبدیل هر جمع به نماینده‌ای هم‌ارز در پیمانهٔ ۱۱، شبکهٔ اصلاح از ثابت
	\(0101_2=5\) استفاده می‌کند: افزودن ۵ در حساب چهاربیتی، همراه با دورریختن نقلی، همان تبدیل \(z\mapsto z-11\) برای مقدارهای نیازمند اصلاح است، زیرا \(z+5-16=z-11\). بیت نقلی جمع اوّلیّه مقدارهای ۱۶ تا ۱۸ را مشخّص می‌کند و مقایسه با مقدار مرزی \(1010_2=10\) نیز مقدارهای ۱۱ تا ۱۵ را از مقدارهای بدون نیاز به اصلاح جدا می‌سازد. این سیگنال‌های کنترلی، ورودی مناسب تسهیم‌کننده‌ها را انتخاب می‌کنند تا حالت‌های \(X=Y\)، \(X=Y+11\) و \(Y=X+11\) به برابری دو مقدار اصلاح‌شده تبدیل شوند. مقایسه‌کنندهٔ نهایی فقط برابری این دو مقدار را می‌سنجد و خروجی برابری آن، خروجی داخلی \lr{Div11} است.

	\shot{screenshots/Div11/6.png}{طرح‌وارهٔ سلسله‌مراتبی آشکارساز مضرب ۱۱}

	\subsection*{ج) مدار نهایی}
	دو آشکارساز مستقل در پیمانهٔ \lr{final} قرار گرفتند و چهار رقم دهدهیِ دودویی‌شدهٔ مشترک به هر دو اعمال شد. نام خروجی‌های داخلی دو زیرمدار به‌ترتیب \lr{div3} و \lr{Div11} است و در سطح نهایی به خروجی‌های \lr{is\_Div3} و \lr{is\_Div11} متّصل می‌شوند. این دو خروجی می‌توانند هم‌زمان و بدون اثرگذاری بر یکدیگر وضعیّت بخش‌پذیری عدد ورودی را گزارش کنند.

	\shot{screenshots/Final/1.png}{اتّصال آشکارسازهای مضرب ۳ و مضرب ۱۱ در مدار نهایی}

	\section*{۱.۳ شبیه‌سازی و ارزیابی}
	در شبیه‌سازی \lr{Div3} چند عدد چهاررقمی اعمال شده است. برای نمونه، ارقام \(1,2,3,6\) عدد ۱۲۳۶ را می‌سازند و چون جمع رقم‌ها \(12\) است، خروجی \lr{div3} یک می‌شود. برای ۲۳۴۱ با جمع رقم‌های \(10\) خروجی صفر، برای ۳۱۰۵ با جمع رقم‌های \(9\) خروجی یک و برای ۵۰۳۸ با جمع رقم‌های \(16\) خروجی دوباره صفر است. این رفتار با معیار نظری بخش‌پذیری بر ۳ سازگار است.

	\shot{screenshots/Div3/9.png}{شبیه‌سازی آشکارساز مضرب ۳ برای چند عدد چهاررقمی}

	شکل موج \lr{Div11} حالت‌های گوناگون دو جمع متناوب را می‌آزماید. برای ۰۰۰۰، دو جمع برابر صفرند و خروجی یک است؛ برای ۱۲۳۲ نیز \(A+C=B+D=4\) و خروجی یک می‌شود. بردار ۱۲۳۴، با دو جمع ۴ و ۶، نمونهٔ اختلاف نامعتبر و دارای خروجی صفر است. بردار ۹۰۲۰ حالت \(X=Y+11\) را می‌آزماید. برای ۹۷۹۰ داریم \(X=9+9=18\) و \(Y=7+0=7\)، پس \(X=Y+11\)؛ برای ۷۹۰۹ داریم \(X=7+0=7\) و \(Y=9+9=18\)، پس \(Y=X+11\). بردار ۹۹۹۹ نیز حالت \(X=Y\) را می‌آزماید. در همهٔ این حالت‌ها خروجی یک است. در پایان، ۹۹۹۸ با اختلاف یک، خروجی صفر تولید می‌کند. بدین‌ترتیب، برابری، هر دو جهت اختلاف ۱۱ و یک اختلاف نامعتبر پوشش داده شده‌اند.

	\shot{screenshots/Div11/7.png}{شبیه‌سازی آشکارساز مضرب ۱۱ برای بردارهای آزمون}

	در آزمون مدار نهایی، هر چهار رقم به هر دو زیرمدار داده شده‌اند. شکل موج نشان می‌دهد که دو خروجی مستقل‌اند: برای ۳۶۶۳ هر دو شرط برقرارند؛ برای ۱۱۱۲ هیچ‌کدام برقرار نیست؛ برای ۵۷۶۳ فقط شرط مضرب ۳ و برای ۶۱۲۷ فقط شرط مضرب ۱۱ برقرار است. این چهار حالت، ترکیب‌های ممکن دو خروجی را پوشش می‌دهند.

	\shot{screenshots/Final/3.png}{شبیه‌سازی یکپارچه و پوشش چهار حالت خروجی دو آشکارساز}

	\section*{۱.۴ سنتز و پیاده‌سازی نرم‌افزاری}
	نمای سطح انتقال ثبات‌های مدار مضرب ۳، مسیر مورد انتظار برای محاسبهٔ باقی‌ماندهٔ هر رقم، جمع نتایج و آشکارسازی باقی‌ماندهٔ صفر را حفظ کرده است. نماهای پس از نگاشت و پس از جانمایی نشان می‌دهند که ساختار دروازه‌ای به منابع واقعی قطعه نگاشت و مسیریابی شده است.

	\shot{screenshots/Div3/RTL_Viewer.png}{نمای \lr{RTL Viewer} آشکارساز مضرب ۳}
	\shot{screenshots/Div3/post_mapping.png}{نمای \lr{Technology Map Viewer (Post-Mapping)} آشکارساز مضرب ۳}
	\shot{screenshots/Div3/post_fitting.png}{نمای \lr{Technology Map Viewer (Post-Fitting)} آشکارساز مضرب ۳}

	در مدار مضرب ۱۱ نیز نمای سطح انتقال ثبات‌ها زنجیرهٔ جمع‌کننده‌ها، تسهیم‌کننده‌ها و مقایسه‌کنندهٔ نهایی را نشان می‌دهد. نماهای فناوری پس از نگاشت و جانمایی، پیاده‌سازی همین مسیر ترکیبی را روی منابع تراشه تأیید می‌کنند.

	\shot{screenshots/Div11/RTL_Viewer.png}{نمای \lr{RTL Viewer} آشکارساز مضرب ۱۱}
	\shot{screenshots/Div11/post_mapping.png}{نمای \lr{Technology Map Viewer (Post-Mapping)} آشکارساز مضرب ۱۱}
	\shot{screenshots/Div11/post_fitting.png}{نمای \lr{Technology Map Viewer (Post-Fitting)} آشکارساز مضرب ۱۱}

	نمای سطح انتقال ثبات‌های مدار نهایی دو پیمانهٔ \lr{Div3} و \lr{Div11} را با ورودی‌های مشترک نشان می‌دهد. در نماهای پس از نگاشت و پس از جانمایی نیز ۱۶ پایهٔ ورودی مربوط به چهار رقم و دو پایهٔ خروجی دیده می‌شوند.

	\shot{screenshots/Final/RTL_Viewer.png}{نمای \lr{RTL Viewer} مدار نهایی}
	\shot{screenshots/Final/post_mapping.png}{نمای \lr{Technology Map Viewer (Post-Mapping)} مدار نهایی}
	\shot{screenshots/Final/post_fitting.png}{نمای \lr{Technology Map Viewer (Post-Fitting)} مدار نهایی}

	هر سه طرح برای قطعهٔ \lr{EP4CGX15BF14C6} از خانوادهٔ \lr{Cyclone IV GX} با موفّقیّت همگردانی شده‌اند. آشکارساز مضرب ۳، تعداد ۳۸ تابع ترکیبی و ۱۷ پایه، آشکارساز مضرب ۱۱، تعداد ۵۴ تابع ترکیبی و ۱۷ پایه، و مدار یکپارچه ۹۲ تابع ترکیبی و ۱۸ پایه مصرف می‌کند. هیچ ثبات، حافظه، ضرب‌کننده یا حلقهٔ قفل‌شده در فاز استفاده نشده و مصرف عناصر منطقی هر سه طرح کمتر از یک درصد ظرفیّت تراشه است؛ این نتیجه با ماهیّت کاملاً ترکیبی مدارها تطابق دارد.

	\shot{screenshots/Div3/10.png}{خلاصهٔ منابع و همگردانی موفّق آشکارساز مضرب ۳}
	\shot{screenshots/Div11/8.png}{خلاصهٔ منابع و همگردانی موفّق آشکارساز مضرب ۱۱}
	\shot{screenshots/Final/2.png}{خلاصهٔ منابع و همگردانی موفّق مدار نهایی}

	شواهد موجود مربوط به شبیه‌سازی و پیاده‌سازی نرم‌افزاری هستند و تصویری از پیکربندی تراشه یا آزمون روی برد ارائه نشده است؛ بنابراین نتیجه‌گیری این گزارش به همین مراحل محدود می‌شود.

	\section*{۱.۵ نتایج}
	\begin{itemize}
		\item زیرمدارهای نیم‌جمع‌کننده، جمع‌کنندهٔ کامل، جمع‌کنندهٔ چهاربیتی، مقایسه‌کننده‌ها و تسهیم‌کننده تنها با طرح‌واره و دروازه‌های پایه ساخته شدند.
		\item معیار جمع رقم‌ها برای تشخیص مضرب ۳ به‌صورت سلسله‌مراتبی پیاده‌سازی و در شبیه‌سازی تأیید شد.
		\item معیار اختلاف جمع رقم‌های متناوب برای تشخیص مضرب ۱۱ با جمع‌کننده، تسهیم‌کننده و مقایسه‌کننده پیاده‌سازی شد.
		\item شبیه‌سازی مدار نهایی هر چهار ترکیب ممکن خروجی‌های \lr{is\_Div3} و \lr{is\_Div11} را نشان داد.
		\item هر سه طرح بدون استفاده از ثبات و با مصرف کمتر از یک درصد عناصر منطقی قطعه سنتز شدند.
	\end{itemize}

	\section*{نتیجه‌گیری}
	در این آزمایش دو مدار ترکیبی تشخیص بخش‌پذیری برای ورودی چهاررقمی با رمزگذاری دهدهیِ دودویی‌شده، به روش طراحی طرح‌واره‌ای ساخته شدند. شکستن طرح به پیمانه‌های کوچک و قابل‌آزمون، از نیم‌جمع‌کننده و مقایسه‌کنندهٔ یک‌بیتی تا آشکارسازهای کامل، بررسی مرحله‌به‌مرحلهٔ عملکرد را ممکن کرد. شکل موج‌ها با قواعد ریاضی بخش‌پذیری بر ۳ و ۱۱ سازگارند و گزارش‌های سنتز و نماهای پس از نگاشت و جانمایی، سنتز و نگاشت موفّق و کم‌هزینهٔ طرح برای قطعهٔ هدف \lr{Cyclone IV GX} را تأیید می‌کنند؛ آزمون فیزیکی روی برد انجام نشده است.
\end{document}
